\section{Checking the scope of privacy claims}
\label{sec:scopechecks}

\subsection{Attack success under broader access}

Adversary privileges govern when empirical attack outcomes transfer across
threat models. A simple monotonic relationship separates demonstrated
breaches from empirical non-detections.

\begin{samepage}
\begin{proposition}[Transfer under nested observations]
\label{prop:nested}
Fix a target distribution, recovery metric, and computational budget.
Suppose observer $B$ receives every value visible to observer $A$ and can
apply every budget-admissible attack for $A$ after projecting onto $A$'s
observation space. There then exists a budget-admissible attack for $B$
whose expected recovery matches or exceeds that of any individual attack
available to $A$. Non-detection by a tested attack under $A$ places no
non-trivial upper bound on recovery under $B$.
\end{proposition}
\begin{proof}
Projecting $B$'s input onto $A$'s view and running $A$'s attack reproduces
the exact output distribution and recovery score within budget. Existence
of that projected attack leaves open whether a particular gradient
optimiser finds it, while failure of one tested attack under $A$ leaves
stronger uses of $B$'s extra state unconstrained.
\end{proof}
\end{samepage}

Shared-key orthogonal wrapping illustrates why
Proposition~\ref{prop:nested} invalidates careless threat-model upgrades
(\artifact{E8}). Multiplying hidden states by an orthogonal matrix $Q$
transforms boundary traffic via
\begin{equation}
 h'=hQ,\qquad h=h'Q^{\mathsf T}.
 \label{eq:rotation}
\end{equation}
Who holds $Q$ defines the security boundary. While a passive network tap
lacking $Q$ sees rotated coordinates $h'$, an authorised cloud server that
stores or derives $Q$ in host RAM to compute $h=h'Q^{\mathsf T}$ breaks the
no-key premise immediately if compromised. Archived experiment accounts and
collection manifests restrict the restoration result to passive wire and
capture observers (\artifact{E8}, \artifact{ER}). Recent attacks that
invert statistical and permutation obfuscation in third-party LLM serving
highlight the same host-exposure boundary \cite{permattack2025}.

Orthogonal maps also preserve Euclidean norms and token inner products
exactly. Coordinate mismatch explains why a position-wise MLP decoder fails
on $h'$ while relational Gram geometry stays intact across the sequence.
Stronger known-plaintext alignment solvers inside the robustness archive
exploit labelled pairs to recover the subspace at depth-dependent sample
sizes (\artifact{E8}), consistent with independent findings that deeper
transformer cuts retain recoverable internal states
\cite{depth2025,actinv2026,bettisplit2026}.

\begin{table}[t]
\centering\small
\begin{tabular}{@{}rrrr@{}}
\toprule
Depth & Estimated $K_{50}$ & Measured bracket & $K_{50}/5120$ \\
\midrule
1 & 258.32 & 200 to 300 & 0.05045 \\
4 & 355.81 & 300 to 400 & 0.06949 \\
8 & 617.93 & 600 to 700 & 0.12069 \\
\bottomrule
\end{tabular}
\caption{Interpolated absolute 50\% token-recovery crossings ($K_{50}$)
from the 27B ($H=5120$) robustness artifact across split depths
(\artifact{E8}). Normalised ratios $K_{50}/5120$ span a $2.39$-fold range
across depth without confidence intervals on the interpolated crossings.}
\label{tab:k50}
\end{table}

Aggregating sessions can obscure peak exposure as well (\artifact{E8}).
Under per-session rotation at depth~1, the pooled top-1 recovery score is
$3.71\%$, yet constituent sessions reach $5.73\%$, $3.65\%$, and $2.34\%$.
Pooled averages understate the leakiest session. Because known-plaintext
solvers extract sensitive partial spans well below the absolute 50\%
crossing in Table~\ref{tab:k50}, interpolated $K_{50}$ brackets track
empirical attack scaling in the lab without defining safe operational key
lifetimes.

\subsection{Comparing rotation cadences}

Archived deployment logs separate a three-session Mistral-7B per-session
rotation test from a later eight-run Mistral-NeMo-12B per-epoch ratcheting
matrix (\artifact{ER}). Per-session swaps add negligible latency. Within
the 12B per-epoch matrix, greedy-decode identity holds in 20 of 24 sessions
and four of the eight runs carry documented admissibility caveats, while
the \texttt{N=256} arm executes zero rotations across its 60-token decode
horizon and acts as a static control (though its zero-RTT throughput still
drops from $7.6$--$7.7$ to $6.3$--$6.4$ tokens/s due to harness overhead).
Wire decoders scored zero top-1 recovery on rotated epochs within that
bounded capture setup.

Frequent ratcheting imposes steep latency penalties (\artifact{ER}). At
cadence \texttt{N=16} with zero emulated RTT in session~1, recorded
throughput drops from $7.6$ tokens/s on the static baseline to $2.8$
tokens/s under ratcheting (a $63.16\%$ reduction at saved precision),
whereas the evidence-triggered budget arm runs near $5.5$ to $5.6$ tokens/s
at zero emulated RTT (and $2.6$ to $2.9$ tokens/s at 80\,ms RTT). Model
scale, rotation cadence, and RTT twins remain distinct experimental factors
across those extracted logs.

Provenance reconciliation explains why two internal documents disagreed
about this matrix (\artifact{ER}). An early summary flagged the entire
ratchet sweep invalid before a subsequent RTT-twin review readmitted the
supporting runs alongside explicit identity, throughput, and cloud-RAM key
caveats. Following our precedence rule, we adopt the later collection
manifest together with every attached operational restriction.

\subsection{Reading mutual information estimates}

Companion boundary evaluations report near-zero Mutual Information Neural
Estimation (MINE) \cite{belghazi2018} readouts on defended latent splits
(\artifact{MI}). For any measurable critic $T$ with finite expectations
($\mathbb E_{P_{XZ}}[|T|]<\infty$ and $\mathbb E_{P_XP_Z}[e^{T}]<\infty$),
the population Donsker--Varadhan variational objective obeys
\begin{equation}
 L_T=\mathbb E_{P_{XZ}}[T(X,Z)]-
 \log\mathbb E_{P_XP_Z}[e^{T(X,Z)}]\le \operatorname{MI}(X,Z),
 \label{eq:dv}
\end{equation}
where $\log$ is the natural logarithm and $\operatorname{MI}(X,Z)$ denotes
true mutual information in nats. That inequality bounds leakage strictly
from below. In addition to severe finite-sample bias and exponential
variance in variational estimators \cite{mcallester2020,song2020},
Eq.~\eqref{eq:dv} points in the wrong direction to certify an upper leakage
budget, just as weak empirical attacks in differential-privacy auditing
never prove small $\varepsilon$ \cite{jagielski2020}.

A trivial analytic channel exposes the flaw immediately. Let input $X$ be
uniform over $m$ discrete symbols, transmit plaintext $Z=X$, and evaluate
constant critic $T=0$. Critic output $L_T$ equals zero identically even
though $\operatorname{MI}(X,Z)=\log m$ nats ($8$ full bits at $m=256$).
Pairing the MINE readout with a second lower-bound estimator (such as the
companion study's calibrated cross-entropy probe bound) leaves the
inequality direction unchanged, and the companion draft's own limitations
section acknowledges that neither estimator furnishes an upper bound
(\artifact{MI}).

We therefore treat near-zero critic readouts as finite non-detections by a
particular neural family while leaving true channel leakage open in both
directions. Architectural ablations require identical restraint
(\artifact{MI}). Observing comparable task metrics across a handful of
invariant MLP and radial remote modules demonstrates empirical saturation
within the tested training budget, whereas proving a universal
information-theoretic capacity ceiling would require bounding every
admissible remote function family.
